We study learned context selection for clinical question answering over Electronic Health Records (EHRs) under strict context budgets. Rather than modifying a downstream language model, we learn a lightweight and interpretable selection layer that scores candidate evidence chunks (note sections and time-stamped structured events) for relevance to a question. Our goal is to empower a smaller, resource-constrained model to improve performance comparable to a larger, more expensive model in the healthcare space. Using MIMIC-IV patient timelines constructed via BigQuery extraction and question--answer supervision generated by a larger model, we train an L1-regularized logistic regression ranker over lexical, semantic, temporal, and section-structure features. The learned selector is then used to augment frozen small LLMs by providing a compact, high-signal context. Our final evaluation protocol compares baseline retrieval methods (BM25/semantic RAG/recency/full-context) against the learned selector, including multi-model generalization and clinician-reviewed benchmarks when available.